\documentclass[11pt,a4paper]{article}

\usepackage{amsmath,amssymb,amsthm}
\usepackage[margin=2.5cm]{geometry}
\usepackage{xcolor}
\usepackage{booktabs}
\usepackage{tabularx}
\usepackage{adjustbox}
\usepackage{enumitem}
\usepackage{microtype}

\usepackage[
  colorlinks=true,
  linkcolor=blue!55!black,
  citecolor=darkgray,
  urlcolor=blue!55!black,
  bookmarks=true,
  bookmarksnumbered=true,
  unicode=true
]{hyperref}

\geometry{a4paper, top=2.5cm, bottom=2.5cm, left=2.5cm, right=2.5cm}

\widowpenalty=10000
\clubpenalty=10000

\theoremstyle{plain}
\newtheorem{theorem}{Theorem}[section]
\newtheorem{lemma}[theorem]{Lemma}
\newtheorem{corollary}[theorem]{Corollary}
\newtheorem{proposition}[theorem]{Proposition}
\theoremstyle{definition}
\newtheorem{definition}[theorem]{Definition}
\newtheorem{example}[theorem]{Example}
\theoremstyle{remark}
\newtheorem{remark}[theorem]{Remark}

\newcommand{\norm}[1]{\lVert #1 \rVert}
\newcommand{\gap}{\operatorname{gap}}
\newcommand{\dist}{\operatorname{dist}}
\newcommand{\R}{\mathbb{R}}
\newcommand{\Z}{\mathbb{Z}}

\hypersetup{
  pdftitle={The Two Diagnostics: Residue Proof Architecture at k=3 and the Deficit Law at k=8,9},
  pdfauthor={MIKEAA2020},
  pdfsubject={Lonely Runner Conjecture, deficit law, bridge attribution, orbit lemma, top-pair rescue},
  pdfkeywords={lonely runner, deficit law, orbit lemma, divisor lemma, top-pair rescue, bridge attribution}
}

\begin{document}

\tableofcontents
\newpage

\section{Introduction}\label{sec:intro}

For a finite set $V$ of $n$ distinct positive integer speeds, the loneliness of $V$ is
\[
  \gap(V) \;=\; \max_{t \in [0,1)} \; \min_{v \in V} \norm{v t},
\]
where $\norm{x}$ denotes the distance from $x$ to the nearest integer. The Lonely Runner
Conjecture (LRC) asserts $\gap(V) \ge 1/(n+1)$ for every $V$. In the preceding audit
(``the C-half attack and exhaustiveness''), the program reduced LRC, conditionally on a
decomposition exhaustiveness statement, to the \emph{deficit law}: at a reduced argmax
$t^{*} = a/b$ of $f_V(t) = \min_v \norm{v t}$, with $m = b \cdot \gap(V)$, the integer
\[
  d \;=\; m(n+1) - b
\]
satisfies $d \ge 0$. The deficit law is equivalent to $\gap(V) \ge 1/(n+1)$; it is LRC in
residue form. Its empirical status was: verified with zero violations on $20{,}804 + 297$
crossings of the drop experiment, not proved; the covering argument of the previous turn
reaches only $b \le n(2m+1)+1$, a factor-$2$ short of the target $b \le m(n+1)$.

The advisor's assessment of that turn isolated the value of the program in the
\emph{frame} --- the clean statement of the factor-$2$ obstruction --- and prescribed two
diagnostics, which this note executes completely.

\begin{enumerate}[label=\textbf{D\arabic*.}, leftmargin=*]
\item \textbf{Independent residue proof at $k = 3$.} Attempt to prove the deficit law for
four speeds by residue methods, without importing the classical LRC proof. If the proof
is genuinely shorter or more generalizable, the reformulation has content; if it merely
re-derives the classical argument, the reformulation is a translation and the factor-$2$
obstruction is the only new thing.
\item \textbf{Counterexample hunt at $k = 8, 9$.} Search for violations of the deficit
law at nine and ten speeds, where LRC is known (Rosenfeld and Trakulthongchai) but not by
residue methods. A failure would be a major finding (bug or counterexample); exceptions,
if any, would be the interesting objects.
\end{enumerate}

Between the two diagnostics we interpose a new experiment, the \emph{bridge attribution}
(Section~\ref{sec:bridge}), which measures --- for every $n = 3, \dots, 10$ --- how much
of the would-be-violation region is closed by the $b$-grid residue arithmetic alone, and
how much requires the off-grid crossing structure. This measurement is what turns the
``clothing versus body'' question into a number.

All computations use exact rational arithmetic. Every would-be violation flagged by the
fast engine is re-verified with the repository's independent reference solver
(\texttt{lrc\_gap\_lib.gap\_int}, the Lemma-6.1 candidate implementation) before being
reported. Reproduction details are collected in Section~\ref{sec:repro}.

\section{The pair-sum candidate theorem}\label{sec:machinery}

The engine behind both diagnostics rests on a sharpened form of the candidate theorem.

\begin{theorem}[pair-sum candidates]\label{thm:pairsum}
Let $V$ be a set of $n \ge 2$ distinct positive speeds. The maximum of
$f_V(t) = \min_{v \in V} \norm{v t}$ over $[0,1)$ is attained at some $t_0 = j/(u+w)$
with $u, w \in V$ and $0 \le j < u + w$.
\end{theorem}

\begin{proof}
The maximum is positive and attained in the open interval $(0,1)$: $f(0) = 0$, and $f$
is strictly positive at, e.g., an odd multiple of $1/(2v)$ minus an epsilon avoiding
the finitely many zeros of the other speeds. Let $t_0 = a/b$ (in lowest terms) be a
maximizing point and $A = \{v \in V : \norm{v t_0} = f_V(t_0)\}$ the active set. In a
neighborhood of $t_0$ the inactive speeds stay strictly above $f_V(t_0)$, so locally
$f_V = \min_{v \in A} \norm{v t}$.

Suppose first that some $v \in A$ sits at its pole, i.e.\ $v t_0 \equiv 1/2 \pmod 1$.
Then $f_V(t_0) = 1/2$ forces $\norm{v' t_0} \ge 1/2$, hence $= 1/2$, for every
$v' \in V$; thus $(v + v') t_0 \equiv 0 \pmod 1$ for every pair, and $t_0$ lies on the
$(v+v')$-grid.

Otherwise all members of $A$ are smooth at $t_0$ with signs determined by their residues
$v a \bmod b$: residue in $(0, b/2)$ gives slope $+v$, residue in $(b/2, b)$ gives slope
$-v$. An interior local maximum cannot have all active slopes of one sign (all $+$
admits an increase to the right; all $-$ admits an increase to the left). Hence $A$
contains $u$ with residue $m_0 \in (0, b/2)$ and $w$ with residue $b - m_0$, where
$m_0 = b \cdot f_V(t_0)$. Then
\[
  (u + w) a \;\equiv\; m_0 + (b - m_0) \;\equiv\; 0 \pmod b .
\]
Since $\gcd(a, b) = 1$, we get $b \mid u + w$, so $t_0 = a/b = j/(u+w)$ with
$j = a (u+w)/b$.
\end{proof}

\begin{remark}
The peak candidates $\mathrm{odd}/(2v)$ and the difference-crossings $j/(w - u)$ of the
Lemma-6.1 candidate list never carry a strict maximum alone: peaks force the all-poles
case, which is on a pair-sum grid, and difference-crossings pair equal residues, i.e.\
equal slopes, which cannot bound a local maximum from both sides. The theorem was
validated empirically: on $2{,}500$ random sets the pair-sum-only engine agrees exactly
with the reference implementation, and every reference argmax time satisfies the
pair-sum property.
\end{remark}

The engine (\texttt{GapEngine}) precomputes the distance tensor
$T[j, D, v] = \min(vj \bmod D,\; D - vj \bmod D)$ over all pair sums $D \le 2M + 1$ and
reduces each evaluation to a gather and two reductions; it processes $10^5$--$10^6$
exact gap evaluations per minute on two cores.

\section{Diagnostic 2: the deficit law at nine and ten speeds}\label{sec:t2}

\subsection{Exhaustive scans}

Every $n$-subset of $[1..M]$ was scanned with the exact engine. Table~\ref{tab:scan}
summarizes the outcome.

\begin{table}[ht]
\centering\small
\begin{tabular}{rrrrrrl}
\toprule
$n$ & $M$ & sets & $\min \gap$ & violations & $d_{\min}$ & critical sets (normalized) \\
\midrule
3 & 18 & 816 & $1/4$ & 0 & 0 & 1 \\
4 & 18 & 3{,}060 & $1/5$ & 0 & 0 & 2 \\
5 & 18 & 8{,}568 & $1/6$ & 0 & 0 & 2 \\
6 & 18 & 18{,}564 & $1/7$ & 0 & 0 & 1 \\
7 & 18 & 31{,}824 & $1/8$ & 0 & 0 & 3 \\
8 & 18 & 43{,}758 & $1/9$ & 0 & 0 & 1 \\
9 & 20 & 167{,}960 & $1/10$ & 0 & 0 & 1 \\
10 & 21 & 352{,}716 & $1/11$ & 0 & 0 & 1 \\
\bottomrule
\end{tabular}
\caption{Exhaustive deficit-law scans: $627{,}266$ sets, zero violations. ``Critical
sets'' are gcd-normalized sets achieving $\gap = 1/(n+1)$ exactly. At $n = 4, 5, 7$
they are $\{1,\dots,n\}$ plus $\{1,3,4,7\}$, $\{1,3,4,5,9\}$,
$\{1,2,3,4,5,7,12\}$, $\{1,4,5,6,7,11,13\}$.}
\label{tab:scan}
\end{table}

Two structural laws of the previous turn were re-verified at the argmax of every single
set: the \emph{pole law} (the residues $v_i a \bmod b$ attain both $m$ and $b - m$) and
the \emph{covering wall} $b \le n(2m+1) + 1$: zero failures across all $627{,}266$ sets,
including all $520{,}676$ sets at nine and ten speeds.

\subsection{Adversarial search and critical isolation}

A gap-minimizing random-restart hill-climb with plateau walks ($M \le 60$) performed
$102{,}122$ evaluations at $n = 9$ and $53{,}260$ at $n = 10$: the search bottoms out at
exactly $1/10$ and $1/11$ respectively, with zero violations.

The critical census deserves its own statement. Non-regular critical sets (gap exactly
$1/(n+1)$, normalized set $\ne \{1, \dots, n\}$) exist at $n = 4, 5, 7$ inside $M = 18$;
each is a floor-landing extension of a rung-2 or tight core, the mechanism the drop
experiment called a tight extension. A targeted hunt at $n = 8, 9, 10$ (growing all
critical and rung-2 seeds at $n-1$ and $n-2$ by one or two fillers, $M \le 60$) found
\emph{no} non-regular critical set; the repository's tight-8 census at $V \le 50$
independently contains exactly the six regular multiples $c\{1,\dots,8\}$. The floor at
$k = 8, 9$ is isolated to the regular family.

\begin{remark}[isolation]
The thinning of the critical census with $n$ is the $k = 8, 9$ analogue of the
``exception'' question in D2: the deficit law holds with no exceptions, and the class
of extremal configurations shrinks to a single point. Any proof of the deficit law at
these sizes must, in residue terms, handle only the regular configuration and the
shallow bridge above it (Section~\ref{sec:bridge}).
\end{remark}

\section{The bridge attribution: clothing versus body, measured}\label{sec:bridge}

The \emph{bridge region} for size $n$ is
\[
  m(n+1) \;<\; b \;\le\; n(2m+1) + 1 :
\]
below it the deficit law holds trivially for a configuration; above it the covering
union bound kills trivially. The bridge is where the factor-$2$ obstruction lives. A
\emph{bridge configuration} is a residue set
$R = \{m\} \cup \{\text{interiors}\} \cup \{b - m\}$ at time $a/b$, with distinct
interiors in $(m, b - m)$. Each configuration is tested at three levels:
\begin{itemize}[leftmargin=*]
\item \textbf{L2 (orbit / $b$-grid):} some $j \in [1, b-1]$ has
$\min_i \dist(j r_i, 0) > m$ --- a violation of the necessary condition of Lemma~B
below;
\item \textbf{L3 (off-grid crossings):} for every unit $s \bmod b$, the canonical lift
$c_i = s\,r_i \bmod b \in [1, b-1]$, viewed as a speed set, has $\gap > m/b$;
\item \textbf{realized:} some lift has $\gap = m/b$ exactly --- an actual violation of
the deficit law.
\end{itemize}

Table~\ref{tab:bridge} gives the kill profile across $n = 3, \dots, 10$ ($m \in \{1,2,3\}$
for $n \ge 5$, $m \in \{1,\dots,8\}$ at $n = 4$; complete enumeration up to 500
configurations per $(m, b)$, seeded sampling beyond; all units for $n \le 7$, 12 sampled
units for $n \ge 8$; $b \le 80$).

\begin{table}[ht]
\centering\small
\begin{tabular}{rrrrr}
\toprule
$n$ & configs & L2 orbit kills & L3 off-grid kills & realized \\
\midrule
3 & 192 & 186 \; (96.9\%) & 6 & 0 \\
4 & 45{,}144 & 45{,}035 \; (99.8\%) & 109 & 0 \\
5 & 14{,}075 & 13{,}781 \; (97.9\%) & 294 & 0 \\
6 & 22{,}321 & 21{,}535 \; (96.5\%) & 786 & 0 \\
7 & 27{,}923 & 26{,}611 \; (95.3\%) & 1{,}312 & 0 \\
8 & 10{,}019 & 9{,}448 \; (94.3\%) & 571 & 0 \\
9 & 11{,}414 & 10{,}725 \; (94.0\%) & 689 & 0 \\
10 & 12{,}804 & 12{,}073 \; (94.3\%) & 731 & 0 \\
\bottomrule
\end{tabular}
\caption{Bridge attribution: $143{,}892$ would-be-violation configurations. The
$b$-grid necessary condition (the residue ``clothing'') kills a stable $\sim$94--99\%
of the bridge at every size; the residue of the residual concentrates in the shallow
bridge (small $m$, $|d| \le 3$).}
\label{tab:bridge}
\end{table}

Three findings. First, \emph{stability}: the orbit share does not decay with $n$; the
$b$-grid arithmetic keeps its covering power at every size. Second, \emph{L2 anatomy}:
about 55\% of orbit kills are available at non-unit multipliers (the divisor lemma,
Lemma~C below), the rest need unit multipliers. Third, \emph{L3 anatomy}: every killing
crossing sits on a pair whose sum $D = u + w$ is not divisible by $b$ (a pair with
$D = b$ lies wholly on the $b$-grid and is capped by the orbit condition; $D = 2b$ is
impossible for speeds below $b$); winners split between $D < b$ and $b < D < 2b$ roughly
evenly, and winning values typically lie in $(m/b,\, 2m/b)$: the kills are decisive.
The L3 survivors concentrate in the \emph{shallow bridge} --- small $m$ and
$|d| = b - m(n+1) \le 3$ --- which is to say: near-regular residue patterns. The deep
bridge is entirely closed by the orbit condition.

\section{Diagnostic 1: the residue proof architecture at \texorpdfstring{$k=3$}{k=3}}\label{sec:t1}

\subsection{Two lemmas with proofs}

\begin{lemma}[orbit]\label{lem:orbit}
Let $t^{*} = a/b$ be an argmax of $f_V$ with value $m/b$ and residues
$r_i = v_i a \bmod b$. Then for every $j \in \{1, \dots, b-1\}$,
\[
  \min_i \dist\!\big(j\, r_i,\; 0\big) \;\le\; m .
\]
Consequently $\bigcup_i \{\, j : \dist(j r_i, 0) \le m \,\} \supseteq \{1, \dots, b-1\}$
and therefore $b \le n(2m+1) + 1$.
\end{lemma}

\begin{proof}
For every time $t$, $f_V(t) \le \gap(V) = m/b$. Apply this at $t = (j a / b) \bmod 1$:
$\norm{v_i \, j a / b} = \dist(j r_i, 0)/b$, so $\min_i \dist(j r_i, 0)/b \le m/b$.
The second statement is the union bound over the $n$ sets
$\{j : \dist(j r_i, 0) \le m\}$, each of size at most $2m+1$.
\end{proof}

The union-bound corollary is the covering wall of the previous turn; Lemma~\ref{lem:orbit}
itself is strictly stronger, and Table~\ref{tab:bridge} measures by exactly how much:
the wall leaves the entire bridge open, while the pointwise orbit condition closes
$\sim$94\% of it.

\begin{lemma}[divisor]\label{lem:divisor}
Let $d$ be a proper divisor of $b$ with $dm < b$. Then $d$ divides some residue $r_i$.
\end{lemma}

\begin{proof}
Apply Lemma~\ref{lem:orbit} at $j = b/d$. Writing $\rho_i = r_i \bmod d$, we have
$(b/d) r_i \equiv (b/d)\rho_i \pmod b$, so $\dist((b/d) r_i, 0) = (b/d) \cdot
\min(\rho_i, d - \rho_i)$. For this to be at most $m$ we need
$\rho_i \le dm/b$ or $\rho_i \ge d - dm/b$; under $dm/b < 1$ both force $\rho_i = 0$.
\end{proof}

\begin{example}
At $n = 4$, $b = 12$, $m = 1$: the divisors $2, 3, 4, 6$ of $12$ must each divide one of
the residues $\{1, r, s, 11\}$; since $11$ is coprime to $12$, one of $r, s$ must be
divisible by $\operatorname{lcm}(4, 6) = 12$ --- impossible inside $(1, 11)$. The
configuration dies before any crossing is examined.
\end{example}

\begin{lemma}[prime counting]\label{lem:prime}
If $b$ is prime, each residue $r_i$ covers exactly the $2m+1$ multipliers
$r_i^{-1}\{-m, \dots, m\}$; hence $b - 1 \le n(2m+1)$, and equality forces an exact
tiling of $\Z_b^{\times}$ by the reflected inverse pairs $\pm r_i^{-1}$.
\end{lemma}

\subsection{The residual case load, completely enumerated}

At $n = 4$ with $m = 1, \dots, 8$, all bridge values of $b$, and all configurations:
exactly \textbf{116} configurations survive the orbit lemma. All 116 are killed by
off-grid pair-sum crossings of their canonical lifts; none is realized. The $m = 1$
layer, in full, with killing witnesses:
\[
\begin{array}{lllll}
b=6:\ \{1,2,3,5\} & t = 1/4 & (u,w) = (1,3), & \text{value } 1/4; &
\{1,3,4,5\} \quad t = 4/9 \\
b=7:\ \{1,2,3,6\} & t = 1/4 & (u,w) = (1,3), & \text{value } 1/4; &
\{1,2,4,6\} \quad t = 3/8 \\
b=7:\ \{1,3,5,6\} & t = 5/11 & (u,w) = (5,6), & \text{value } 3/11; &
\{1,4,5,6\} \quad t = 4/9 \\
b=8:\ \{1,3,4,7\} & t = 1/5 & (u,w) = (1,4), & \text{value } 1/5; &
\{1,4,5,7\} \quad t = 1/3
\end{array}
\]
(every value exceeds the target $m/b = 1/6, \dots, 1/8$ by at least a factor
$3/2$). The case count is completed by two structural laws.

\paragraph{The dilation law.}
The 20 survivors at $m = 2$ are all within $+1$ (per gap) of $2\times$-dilations of the
eight $m = 1$ patterns; the 10 survivors at $m = 3$ are all within $+1$ of
$3\times$-dilations; 8 of the 20 at $m = 2$ are exact dilations. Exact dilations carry
identical witnesses (a dilated set is a scaled set, the gap is scale-invariant, and the
killing time divides). The residual case load is therefore \emph{parametric} in $m$:
eight base patterns plus bounded perturbations --- the ``consecutive plus one
modification'' structure guessed in the previous turn's open problem O1 is confirmed and
is the whole story at $k = 3$.

\paragraph{The top-pair rescue law.}
For a bridge configuration surviving the orbit lemma, let $c_1 < \dots < c_n$ be any
unit lift. Claim: the pair-sum crossing family of the two \emph{largest} speeds (the
residues nearest the down-pole $b - m$) contains a time at which the envelope exceeds
$m/b$. Measured coverage over all orbit-surviving (configuration, lift) pairs:

\begin{table}[ht]
\centering\small
\begin{tabular}{lrrrrrrr}
\toprule
$n$ & 4 & 5 & 6 & 7 & 8 & 9 & 10 \\
\midrule
(config, lift) pairs & 1{,}798 & 3{,}006 & 8{,}208 & 15{,}212 & 5{,}280 & 6{,}966 & 7{,}194 \\
top-pair coverage & 92.8\% & 94.9\% & 97.1\% & 97.0\% & 97.7\% & 98.3\% & 98.6\% \\
\bottomrule
\end{tabular}
\caption{Top-pair rescue coverage, rising with $n$. On the canonical lift the top pair
kills $116/116$ residual cases at $n = 4$ and $320/323$ at $n = 5$; the three exceptions
fall to bottom and mixed pairs. Whenever the top pair fails, at least two (typically
three at $n = 4$; five to thirty-five at $n = 10$) other winning pairs exist.}
\label{tab:rescue}
\end{table}

\subsection{Assembly status and the honest comparison}

The residue proof architecture at $k = 3$ is:
\[
\text{deficit law at } n = 4
\;=\; \underbrace{\text{Lemma~\ref{lem:orbit}}}_{2 \text{ lines}}
\;+\; \underbrace{\text{Lemma~\ref{lem:divisor}}}_{3 \text{ lines}}
\;+\; \text{top-pair rescue,}
\]
where the last term is the top-pair rescue for eight base patterns and their bounded
perturbations.
The first two lemmas are proved and $n$-uniform. The rescue is the single remaining gap:
one uniform statement, verified for $m \le 8$ and all unit lifts at $n = 4$, and with
the measured coverage of Table~\ref{tab:rescue} at every size.

Compared with the classical proofs at four speeds (and the Barajas--Serra machinery up
to seven), which are covering/potential arguments with $n$-specific case analysis whose
pigeonhole content is the union bound $b \le n(2m+1)+1$: the residue chain does not
re-derive them. Its two proved lemmas kill $99.8\%$ of the $n = 4$ bridge with no case
analysis at all, and the residual is a parametric family with a uniform witness rule
that has no classical counterpart. But the work is not avoided --- it is relocated into
the single rescue statement, which is precisely the factor-$2$ obstruction in disguise.
The verdict on D1 is therefore \emph{between} ``translation'' and ``content'': the
reformulation yields a genuinely different, $n$-uniform proof \emph{architecture} with
measured covering power, and it localizes the difficulty to a thin, highly structured
set --- but completing the proof requires proving the rescue law.

\section{Verdict and recommendation}\label{sec:verdict}

\begin{itemize}[leftmargin=*]
\item \textbf{D1 (residue proof at $k=3$): not a pure translation, not a proof.} The
orbit lemma is strictly stronger than classical covering (it kills what the union bound
cannot reach); the divisor lemma and the top-pair rescue law are new; the residual is
parametric. The factor-$2$ wall is \emph{localized} --- to the shallow bridge,
$|d| \le 3$, near-regular residue patterns --- not removed.
\item \textbf{D2 (counterexample hunt at $k = 8, 9$): clean hold.} $627{,}266$
exhaustive sets, $155{,}382$ adversarial evaluations, $143{,}892$ bridge
configurations: zero violations, zero exceptions, zero machinery discrepancies. The
floor at $n \ge 8$ is regular-isolated.
\item \textbf{Clothing versus body:} the residue clothing closes $\sim$94\% of the body
at \emph{every} size $n = 3, \dots, 10$; the remaining $\sim$6\% --- the shallow bridge
--- is where LRC actually lives in this language.
\end{itemize}

The advisor's decision rule (both tests neutral $\Rightarrow$ pivot) is not quite
triggered: the tests returned mildly positive but unproved. Two moves are now
well-posed.

\paragraph{R1 (one bounded residue-track mile).} Prove the top-pair rescue lemma: for a
bridge configuration with residues $c_1 < \dots < c_n < b$ (any unit lift), the family
$\{j/(c_{n-1} + c_n)\}$ contains $t$ with $\min_i \norm{c_i t} > m/b$. Success completes
the $k = 3$ proof and gives the architecture for $k = 4, 5$.

\paragraph{R2 (the pivot, sharpened).} The orbit lemma is itself a covering statement in
the group language: \emph{the dilates $jR$ of the residue set always meet the arc
$[-m, m]$}. This is exactly the covering formulation of the kernel dictionary (bridge
prospectus, Step 1). The group-language track and the residue track therefore meet at
Lemma~\ref{lem:orbit}: additive-combinatorial bounds on multiplicative covering of arcs
in $\Z_b$ would simultaneously strengthen the orbit lemma and import the dictionary's
first inequality. The pivot is a continuation by other means, aimed at the shallow
bridge, with Table~\ref{tab:bridge} as the target map.

\section{Reproduction}\label{sec:repro}

All scripts live in \texttt{scripts/} (mirrored to \texttt{computational-check/} \allowbreak
\texttt{diag2/}) with outputs in \texttt{scripts/out/} \allowbreak
\texttt{diag2\_*.json}: \texttt{lrc\_diag\_gap.py} (engine + validation),
\texttt{lrc\_diag\_scan.py} (exhaustive scans), \texttt{lrc\_diag\_climb.py}
(adversarial hill-climb), \texttt{lrc\_diag\_critical\_hunt.py} (critical-set hunt),
\texttt{lrc\_diag\_bridge.py} (bridge attribution), \texttt{lrc\_diag\_casebook.py}
(complete residual case book at $n = 4, 5$), \texttt{lrc\_diag\_rescue.py} (top-pair
rescue coverage). Runtimes: the $n = 10$ scan is $52$~s; everything else runs in
seconds to a minute on two cores; exact rational arithmetic throughout.

\end{document}
